\section{Core primitives}

Let $N=n+1$.  For $u\in\F_q^n$, write $\tilde{u}=(-u,1)\in\F_q^N$.

\subsection{Sparse-LPN zero encryption}

\begin{construction}[Sparse-LPN zero vector encryption]
\label{con:zero-vector}
For secret $u\in\F_q^n$, sample $a\leftarrow S_{n,k,q}$ and $e\leftarrow\chi_{q,\eta}$.  Output
\[
    z=(a,a\transpose u+e)\in\F_q^{N}.
\]
Then
\[
    z\cdot\tilde{u}=e.
\]
The row support is at most $k+1$.
\end{construction}

This is a noisy encryption of zero.  Adding a message in the last coordinate gives compact message encryption.

\subsection{Compact input encryption}

\begin{construction}[Compact sparse-LPN input encryption]
\label{con:input-enc}
For secret $t\in\F_q^n$ and message $\mu\in\F_q$:
\begin{enumerate}
    \item Sample $z=(a,a\transpose t+e)$ as in Construction~\ref{con:zero-vector}.
    \item Output
    \[
        c=z+\mu e_N=(a,a\transpose t+e+\mu),
    \]
    where $e_N=(0,\ldots,0,1)\in\F_q^N$.
\end{enumerate}
Then
\[
    c\cdot\tilde{t}=\mu+e.
\]
\end{construction}

The decryption error is not rounded away; it is corrected statistically by repetition and final plurality.  This is natural in the LPN/code setting.

\subsection{Sparse-LPN GSW-style zero matrix}

\begin{construction}[Sparse-LPN GSW-style zero matrix]
\label{con:gsw-zero}
For secret $s\in\F_q^n$, sample $N$ independent sparse-LPN zero rows
\[
    z_j=(a_j,a_j\transpose s+e_j),\qquad j=1,\ldots,N.
\]
Let $Z\in\F_q^{N\times N}$ be the matrix with rows $z_j$.  Then
\[
    Z\tilde{s}=e,
\]
where $e=(e_1,\ldots,e_N)\transpose$.  Every row has support at most $k+1$.
\end{construction}

\subsection{GSW-style scalar encryption}

\begin{construction}[Sparse-LPN GSW-style scalar encryption]
\label{con:gsw-scalar}
To encrypt $\alpha\in\F_q$ under $s$:
\begin{enumerate}
    \item Sample $Z$ as in Construction~\ref{con:gsw-zero}.
    \item Output
    \[
        X=Z+\alpha I_N.
    \]
\end{enumerate}
Then
\[
    X\tilde{s}=\alpha\tilde{s}+e.
\]
Each row has support at most $k+2$.
\end{construction}

This formulation is important for security: $X$ is a pseudorandom zero-encryption matrix plus a public diagonal shift.  This mirrors the approximate-eigenvector technique in GSW-style encryption \cite{GSW2013} and its sparse-LPN adaptation \cite{CorriganGibbsHenzingerKalaiVaikuntanathan2024}.

\subsection{Expansion from compact input to GSW form}

The expansion key contains GSW encryptions of the coordinates of $\tilde{t}$:
\[
    \mathsf{EK}^{\rm exp}_i=\GSW\Enc_s(\tilde{t}_i),
    \qquad i=1,\ldots,N.
\]
Given compact ciphertext $c=(c_1,\ldots,c_N)$, define
\[
    \Expand(c)=\sum_{i=1}^{N}c_i\mathsf{EK}^{\rm exp}_i.
\]
Since $c$ is sparse with at most $k+1$ nonzero coordinates, expansion sums at most $k+1$ GSW matrices.

\begin{lemma}[Expansion identity]
\label{lem:expansion-identity}
Let $c$ be a compact encryption of $\mu$ under $t$.  Let $C=\Expand(c)$.  Then
\[
    C\tilde{s}=\mu\tilde{s}+e',
\]
where
\[
    e'=e_0\tilde{s}+\sum_{i\in\supp(c)}c_i e^{(i)}.
\]
Here $e_0=c\cdot\tilde{t}-\mu$ is the input error and $e^{(i)}$ is the GSW error vector in $\mathsf{EK}^{\rm exp}_i$.
\end{lemma}

\begin{proof}
For each $i$,
\[
    \mathsf{EK}^{\rm exp}_i\tilde{s}=\tilde{t}_i\tilde{s}+e^{(i)}.
\]
Therefore
\[
    C\tilde{s}
    =\sum_i c_i\tilde{t}_i\tilde{s}+\sum_i c_ie^{(i)}
    =(c\cdot\tilde{t})\tilde{s}+\sum_i c_ie^{(i)}.
\]
Since $c\cdot\tilde{t}=\mu+e_0$, the claim follows.
\end{proof}

\subsection{Code-LPN linearly homomorphic encryption}

Let $\mathsf{Code}:\F_q\to\F_q^{m_c}$ be a linear error-correcting encoding.  For a first prototype, $\mathsf{Code}(x)=x\one$.  Let $\mathsf{Decode}$ be its decoder.

\begin{construction}[CLPN encryption]
\label{con:clpn}
For secret $r\in\F_q^{n_c}$ and message $x\in\F_q$:
\begin{enumerate}
    \item Sample $A\leftarrow\F_q^{m_c\times n_c}$.
    \item Sample $e\in\F_q^{m_c}$ with independent $\chi_{q,\eta_c}$ coordinates.
    \item Output
    \[
        Z=(A,b=Ar+\mathsf{Code}(x)+e).
    \]
\end{enumerate}
Decryption computes
\[
    y=b-Ar=\mathsf{Code}(x)+e
\]
and outputs $\mathsf{Decode}(y)$.
\end{construction}

\begin{construction}[CLPN linear evaluation]
Given ciphertexts $Z_i=(A_i,b_i)$ encrypting $x_i$ and coefficients $\alpha_i\in\F_q$, output
\[
    \sum_i \alpha_i Z_i
    =\left(\sum_i\alpha_iA_i,\sum_i\alpha_i b_i\right).
\]
This decrypts to
\[
    \mathsf{Code}\left(\sum_i\alpha_ix_i\right)+\sum_i\alpha_ie_i.
\]
\end{construction}
